\subsubsection{Manual}\label{trien-khai-remediate-manual}

Khắc phục thủ công được áp dụng đối với các khuyến nghị về thiết kế hạ tầng hoặc yêu cầu can thiệp sâu của quản trị viên hệ thống để tránh làm gián đoạn dịch vụ.

\textbf{Bảng chi tiết các hoạt động Remediation thủ công:}

{\def\LTcaptype{none}
\begin{center}
\captionof{table}{Bảng chi tiết khắc phục thủ công}\label{tab:chi-tiet-khac-phuc-thu-cong}
\end{center}
\begin{longtable}[]{@{}|
  >{\raggedright\arraybackslash}p{(\linewidth - 8\tabcolsep - 5\arrayrulewidth) * \real{0.0739}}|
  >{\raggedright\arraybackslash}p{(\linewidth - 8\tabcolsep - 5\arrayrulewidth) * \real{0.2995}}|
  >{\raggedright\arraybackslash}p{(\linewidth - 8\tabcolsep - 5\arrayrulewidth) * \real{0.3000}}|
  >{\raggedright\arraybackslash}p{(\linewidth - 8\tabcolsep - 5\arrayrulewidth) * \real{0.3266}}|@{}}
\hline
\begin{minipage}[b]{\linewidth}\raggedright
\textbf{Phần}
\end{minipage} & \begin{minipage}[b]{\linewidth}\raggedright
\textbf{Benchmark}
\end{minipage} & \begin{minipage}[b]{\linewidth}\raggedright
\textbf{Phương pháp thực hiện}
\end{minipage} & \begin{minipage}[b]{\linewidth}\raggedright
\textbf{Hành động thực hiện}
\end{minipage} \\
\hline
\endhead

1 & Phân vùng riêng cho các container & Tạo phân vùng LVM riêng và định dạng ext4 & Dừng Docker, backup dữ liệu, mount phân vùng mới, cập nhật /etc/fstab, và khởi động lại Docker \\
\hline
& Chỉ cho trusted user kiểm soát Docker Daemon & Sử dụng gpasswd để xóa các user không được phép khỏi nhóm docker & Chạy lệnh sudo gpasswd -d <username> docker để gỡ bỏ user \\
\hline
& Cấu hình kiểm tra nhật ký cho Docker Daemon & Cấu hình audit rule cho dockerd trong docker.rules & Ghi cấu hình -w /usr/bin/dockerd -k docker và reload auditd \\
\hline
& Audit folder /run/containerd & Cấu hình audit rule cho /run/containerd trong docker.rules & Ghi cấu hình -a exit,always -F path=/run/containerd -F perm=war -k docker và reload auditd \\
\hline
& Audit folder /var/lib/docker & Cấu hình audit rule cho /var/lib/docker trong docker.rules & Ghi cấu hình -a exit,always -F path=/var/lib/docker -F perm=war -k docker và reload auditd \\
\hline
& Audit folder /etc/Docker & Cấu hình audit rule cho /etc/docker trong docker.rules & Ghi cấu hình -w /etc/docker -k docker và reload auditd \\
\hline
& Audit docker.service & Cấu hình audit rule cho docker.service trong docker.rules & \begin{minipage}[t]{\linewidth}\raggedright
Ghi cấu hình -w /usr/lib/systemd/\\
system/docker.service -k docker\strut
\end{minipage} \\
\hline
& Audt containerd.sock & Cấu hình audit rule cho containerd.sock trong docker.rules & Ghi cấu hình -w /run/containerd/containerd.sock -k docker \\
\hline
& Audit docker.sock & Cấu hình audit rule cho docker.sock trong docker.rules & Ghi cấu hình -w /run/docker.sock -k docker \\
\hline
& Audt /etc/default/docker & Cấu hình audit rule cho /etc/default/docker trong docker.rules & Ghi cấu hình -w /etc/default/docker -k docker \\
\hline
& Audit /etc/docker/daemon.json & Cấu hình audit rule cho /etc/docker/daemon.json trong docker.rules & Ghi cấu hình -w /etc/docker/daemon.json -k docker \\
\hline
& Audit /etc/containerd/config.toml & Cấu hình audit rule cho /etc/containerd/config.toml & Ghi cấu hình -w /etc/containerd/config.toml -k docker \\
\hline
& Aduti /etc/sysconfig/docker & Cấu hình audit rule cho /etc/sysconfig/docker & Ghi cấu hình -w /etc/sysconfig/docker -k docker \\
\hline
& Audit /usr/bin/containerd-shim & Cấu hình audit rule cho /usr/bin/containerd-shim & Ghi cấu hình -w /usr/bin/containerd-shim -k docker \\
\hline
& Audit /usr/bin/containerd-shim-runc-v1 & Cấu hình audit rule cho containerd-shim-runc-v1 & Ghi cấu hình -w /usr/bin/containerd-shim-runc-v1 -k docker \\
\hline
& Audit /usr/bin/containerd-shim-runc-v2 & Cấu hình audit rule cho containerd-shim-runc-v2 & Ghi cấu hình -w /usr/bin/containerd-shim-runc-v2 -k docker \\
\hline
& Audit /usr/bin/runc & Cấu hình audit rule cho runc & Ghi cấu hình -w /usr/bin/runc -k docker \\
\hline
& Đảm báo bảo mật cho container & Cấu hình hardening OS của container host & Cập nhật OS, tắt dịch vụ thừa, cấu hình SSH và tường lửa \\
\hline
& Đảm bảo phiên bản mới cho Docker kt & Cập nhật Docker lên phiên bản mới nhất & Cập nhật danh sách gói và chạy apt install --only-upgrade docker-ce \\
\hline

2 & Chạy docker daemon tránh quyền root & Cài đặt và cấu hình Docker rootless mode nếu môi trường hỗ trợ. & Tạo user riêng cho Docker, chạy \texttt{dockerd-rootless-setuptool.sh install}, cấu hình biến môi trường và kiểm tra lại bằng \texttt{docker info}. Nếu môi trường production chưa hỗ trợ, ghi nhận lý do và áp dụng kiểm soát truy cập chặt chẽ cho nhóm \texttt{docker}. \\
\hline
& Network traffic được đặt hạn chế giữa các container ở default bridge & Chỉnh cấu hình Docker daemon để hạn chế inter-container communication. & Cập nhật \texttt{/etc/docker/daemon.json} với tùy chọn \texttt{"icc": false}, sau đó restart Docker và kiểm tra lại network mặc định. \\
\hline
& Logging level set ở ``info'' & Cấu hình mức log của Docker daemon trong \texttt{daemon.json}. & Thêm hoặc sửa \texttt{"log-level": "info"}; không sử dụng \texttt{debug} trong môi trường production vì có thể sinh nhiều log và lộ thông tin nhạy cảm. \\
\hline
& Docker được quyền thay đổi iptables & Đảm bảo Docker không bị vô hiệu hóa quyền quản lý iptables. & Xóa cấu hình \texttt{"iptables": false} nếu không có giải pháp network thay thế, hoặc đặt \texttt{"iptables": true}, sau đó restart Docker. \\
\hline
& Đảm bảo không sử dụng registry không an toàn & Loại bỏ insecure registry khỏi cấu hình Docker daemon. & Xóa mục \texttt{insecure-registries} trong \texttt{daemon.json}; chuyển registry nội bộ sang HTTPS/TLS hoặc chỉ cho phép registry đã được phê duyệt. \\
\hline
& Đảm bảo không sử dụng trình điều khiển aufs & Chuyển storage driver sang driver được khuyến nghị. & Backup image, container và volume cần thiết; cấu hình \texttt{"storage-driver": "overlay2"}; sau đó restart Docker và kiểm tra lại bằng \texttt{docker info}. \\
\hline
& Đảm bảo không dùng devicemapper & Không sử dụng \texttt{devicemapper} ở chế độ không phù hợp cho production. & Lập kế hoạch migrate sang \texttt{overlay2}; backup dữ liệu Docker trước khi đổi driver để tránh mất image/container/volume. \\
\hline
& Đảm bảo xác thực TLS cho Docker Daemon & Cấu hình TLS mutual authentication nếu Docker daemon mở TCP API. & Tạo CA, server certificate và client certificate; cấu hình \texttt{tlsverify}, \texttt{tlscacert}, \texttt{tlscert}, \texttt{tlskey} trong \texttt{daemon.json} hoặc systemd override. \\
\hline
& Đảm bảo giới hạn ulimit mặc định & Thiết lập default ulimit cho Docker daemon. & Thêm \texttt{default-ulimits} vào \texttt{daemon.json}, ví dụ giới hạn \texttt{nofile} hoặc \texttt{nproc}; restart Docker và kiểm tra lại container mới. \\
\hline
& Hỗ trợ user namespace & Bật user namespace remapping nếu ứng dụng tương thích. & Cấu hình \texttt{"userns-remap": "default"} trong \texttt{daemon.json}; kiểm tra volume, permission và ứng dụng trước khi áp dụng production. \\
\hline
& Cgroup sử dụng được xác nhận & Kiểm tra và cấu hình \texttt{cgroup-parent} nếu hệ thống có chính sách tài nguyên riêng. & Nếu không có yêu cầu đặc biệt thì giữ mặc định; nếu có, khai báo \texttt{cgroup-parent} rõ ràng và ghi nhận trong tài liệu vận hành. \\
\hline
& Đảm bảo device size không thay đổi trừ khi cần thiết & Loại bỏ cấu hình base device size không cần thiết. & Xóa tùy chọn \texttt{storage-opt dm.basesize} nếu không có nhu cầu; nếu bắt buộc dùng thì cần có phê duyệt và đánh giá tác động. \\
\hline
& Đảm bảo authorization cho Docker client & Triển khai authorization plugin khi cần kiểm soát quyền Docker API. & Cài đặt và cấu hình plugin phân quyền; khai báo \texttt{authorization-plugins} trong cấu hình Docker daemon, sau đó kiểm thử quyền truy cập. \\
\hline
& Đảm bảo container bị hạn chế không cho quyền mới & Bật chính sách \texttt{no-new-privileges} ở cấp daemon hoặc container. & Thêm \texttt{"no-new-privileges": true} vào \texttt{daemon.json} hoặc cấu hình \texttt{security-opt=no-new-privileges} khi chạy container. \\
\hline
& Đảm bảo container có live restore & Bật tính năng live restore của Docker daemon. & Thêm \texttt{"live-restore": true} vào \texttt{daemon.json}; restart Docker và kiểm tra container vẫn tiếp tục chạy khi daemon restart. \\
\hline
& Đảm bảo có log cho đăng nhập gần và từ xa & Cấu hình logging driver hoặc logging tập trung. & Sử dụng \texttt{json-file} có giới hạn dung lượng, hoặc chuyển sang \texttt{syslog}, \texttt{fluentd}, \texttt{gelf} tùy hệ thống log tập trung. \\
\hline
& Đảm bảo Userland Proxy đã tắt & Tắt userland proxy nếu không cần thiết. & Thêm \texttt{"userland-proxy": false} vào \texttt{daemon.json}; restart Docker và kiểm tra lại các dịch vụ publish port. \\
\hline
& Đảm bảo daemon-wide custom seccomp profile được thiết đặt & Thiết kế và áp dụng seccomp profile phù hợp. & Sử dụng default seccomp profile hoặc custom profile đã kiểm thử; không dùng \texttt{seccomp=unconfined} nếu không có phê duyệt. \\
\hline
& Đảm bảo các tính năng thử nghiệm không được triển khai trong môi trường sản xuất & Tắt experimental features trên Docker daemon. & Đảm bảo \texttt{"experimental": false} hoặc không khai báo experimental trong \texttt{daemon.json}; restart Docker và kiểm tra lại bằng \texttt{docker version}. \\
\hline
3 & Đảm bảo quyền root:root cho docker.service & Sửa owner/group của file \texttt{docker.service}. & Chạy \texttt{sudo chown root:root <FragmentPath>} sau khi xác định đường dẫn bằng \texttt{systemctl show -p FragmentPath docker.service}. \\
\hline
& Đảm bảo docker.service cho quyền hợp lý & Sửa permission của file \texttt{docker.service}. & Chạy \texttt{sudo chmod 644 <FragmentPath>} hoặc quyền hạn chế hơn, sau đó reload systemd nếu cần. \\
\hline
& Đảm bảo quyền sở hữu của docker.socket được đặt thành root:root & Sửa owner/group của file \texttt{docker.socket}. & Xác định đường dẫn bằng \texttt{systemctl show -p FragmentPath docker.socket}, sau đó chạy \texttt{sudo chown root:root <FragmentPath>}. \\
\hline
& Đảm bảo quyền truy cập tệp docker.socket được đặt thành 644 hoặc hạn chế & Sửa permission của file \texttt{docker.socket}. & Chạy \texttt{sudo chmod 644 <FragmentPath>} hoặc quyền hạn chế hơn để tránh sửa đổi trái phép socket unit. \\
\hline
& Đảm bảo quyền sở hữu thư mục /etc/docker được đặt thành root:root & Sửa owner/group của thư mục cấu hình Docker. & Chạy \texttt{sudo chown root:root /etc/docker}; kiểm tra lại bằng \texttt{stat -c \%U:\%G /etc/docker}. \\
\hline
& Đảm bảo quyền truy cập thư mục /etc/docker được đặt thành 755 hoặc hạn chế & Sửa permission thư mục \texttt{/etc/docker}. & Chạy \texttt{sudo chmod 755 /etc/docker} hoặc quyền hạn chế hơn nếu phù hợp. \\
\hline
& Đảm bảo quyền truy cập tệp docker.socket được đặt thành 644 hoặc hạn chế & Sửa permission của Docker socket liên quan. & Kiểm tra \texttt{/var/run/docker.sock} và file unit socket; điều chỉnh quyền theo khuyến nghị, tránh để socket có quyền ghi quá rộng. \\
\hline
& Đảm bảo quyền root:root cho sở hữu tệp chứng chỉ & Sửa owner/group của các file certificate. & Chạy \texttt{sudo chown root:root /etc/docker/certs.d/*} với các file chứng chỉ phù hợp; kiểm tra tránh áp dụng sai cho thư mục rỗng. \\
\hline
& Đảm bảo quyền truy cập chứng chỉ registry được đặt thành 444 hoặc hạn chế hơn & Sửa permission của certificate registry. & Chạy \texttt{sudo chmod 444} cho certificate public; đối với private key dùng quyền chặt hơn như \texttt{400} hoặc \texttt{440}. \\
\hline
& Đảm bảo quyền sở hữu Docker server là root:root và file permission 444 & Sửa owner/group và permission của Docker server certificate. & Chạy \texttt{sudo chown root:root <path\_to\_tlscert>} và \texttt{sudo chmod 444 <path\_to\_tlscert>}; private key phải dùng quyền hạn chế hơn. \\
\hline
& Đảm bảo quyền Docker socket file ownership là root:docker và để file permission 660 & Sửa quyền của \texttt{/var/run/docker.sock}. & Chạy \texttt{sudo chown root:docker /var/run/docker.sock} và \texttt{sudo chmod 660 /var/run/docker.sock}. \\
\hline
& Đảm bảo quyền daemon.json là root:root và file permission 644 & Sửa owner/group và permission của \texttt{daemon.json}. & Chạy \texttt{sudo chown root:root /etc/docker/daemon.json} và \texttt{sudo chmod 644 /etc/docker/daemon.json}. \\
\hline
& Đảm bảo quyền /etc/sysconfig/docker là root:root và file permission 644 & Sửa permission file \texttt{/etc/sysconfig/docker} nếu tồn tại. & Chạy \texttt{sudo chown root:root /etc/sysconfig/docker} và \texttt{sudo chmod 644 /etc/sysconfig/docker}; nếu Ubuntu không có file này thì ghi N/A. \\
\hline
& Đảm bảo containerd socket có quyền root:root và permission 660 & Sửa quyền của containerd socket. & Chạy \texttt{chown root:root} và \texttt{chmod 660} cho containerd socket; kiểm tra lại dịch vụ containerd sau khi đổi quyền. \\
\hline
4 & Phải có container tạo, chạy trong trusted base images và không có package dư thừa không cần thiết được tải về. & Chọn base image chính thống, tối giản và loại bỏ package thừa. & Sửa Dockerfile để dùng image đáng tin cậy như \texttt{ubuntu}, \texttt{debian-slim}, \texttt{alpine} hoặc image nội bộ đã phê duyệt; gỡ package không cần thiết sau khi cài. \\
\hline
& Đảm bảo image được scan và phải có security patches & Quét image và rebuild khi có bản vá bảo mật. & Dùng Trivy/Grype/Clair để scan image; cập nhật base image, rebuild và redeploy image sau khi vá CVE nghiêm trọng. \\
\hline
& Phải có HEALTHCHECK instruction & Thêm \texttt{HEALTHCHECK} vào Dockerfile hoặc cấu hình runtime. & Bổ sung lệnh kiểm tra trạng thái ứng dụng, ví dụ kiểm tra HTTP endpoint hoặc process chính; build lại image và test trạng thái healthy. \\
\hline
& Phải có update instruction và không sử dụng một mình trong Dockerfiles & Gộp lệnh update và install trong cùng một lớp build. & Sửa Dockerfile để dùng một lệnh \texttt{RUN apt-get update \&\& apt-get install ...}; dọn cache package sau khi cài đặt. \\
\hline
& Không cung cấp quyền Setuid và Getuid & Loại bỏ setuid/setgid bit không cần thiết trong image. & Kiểm tra file có setuid/setgid trong image, sau đó sửa Dockerfile để xóa quyền đặc biệt bằng \texttt{chmod a-s} đối với file không cần thiết. \\
\hline
& Dùng COPY thay vì ADD & Thay lệnh \texttt{ADD} bằng \texttt{COPY} khi chỉ cần sao chép file. & Sửa Dockerfile, dùng \texttt{COPY} cho file local; chỉ dùng \texttt{ADD} khi thực sự cần chức năng giải nén hoặc tải từ URL. \\
\hline
& Secrets không được bỏ vào Dockerfiles & Loại bỏ secrets khỏi Dockerfile và image layer. & Xóa password, token, private key khỏi Dockerfile; dùng Docker secrets, Ansible Vault, biến môi trường bảo mật hoặc secret manager. \\
\hline
& Các package phải chính thống khi tải về & Chỉ dùng repository và package có nguồn gốc đáng tin cậy. & Cấu hình repository chính thống, kiểm tra GPG key, pin version nếu cần và tránh tải package từ nguồn không rõ ràng. \\
\hline
& Các signed artifacts phải được kiểm chứng & Xác minh chữ ký artifact trước khi sử dụng hoặc upload. & Dùng GPG, cosign hoặc cơ chế ký nội bộ để verify image/artifact; chỉ đưa artifact đã xác minh vào registry hoặc pipeline. \\
\hline
5 & 5.1: Không bật swarm mode nếu không cần thiết & \begin{minipage}[t]{\linewidth}\raggedright
Rời Swarm nếu host không dùng cluster.\strut
\end{minipage} & \begin{minipage}[t]{\linewidth}\raggedright
Chạy docker swarm leave --force sau khi xác nhận không còn workload Swarm.\strut
\end{minipage} \\
\hline
 & 5.2: Bật AppArmor Profile nếu môi trường hỗ trợ & \begin{minipage}[t]{\linewidth}\raggedright
Khởi chạy container với AppArmor profile phù hợp.\strut
\end{minipage} & \begin{minipage}[t]{\linewidth}\raggedright
Dùng security-opt apparmor profile hoặc profile mặc định đã phê duyệt.\strut
\end{minipage} \\
\hline
 & 5.3: Thiết lập SELinux security options nếu áp dụng & \begin{minipage}[t]{\linewidth}\raggedright
Chỉ áp dụng SELinux label khi host thực sự bật SELinux.\strut
\end{minipage} & \begin{minipage}[t]{\linewidth}\raggedright
Bật SELinux đúng môi trường hoặc bỏ label không có tác dụng trên Ubuntu/AppArmor.\strut
\end{minipage} \\
\hline
 & 5.4: Hạn chế Linux kernel capabilities & \begin{minipage}[t]{\linewidth}\raggedright
Drop capability mặc định và chỉ add lại capability cần thiết.\strut
\end{minipage} & \begin{minipage}[t]{\linewidth}\raggedright
Dùng cap-drop ALL; kiểm thử rồi chỉ cap-add capability đã được phê duyệt.\strut
\end{minipage} \\
\hline
 & 5.5: Không sử dụng privileged container & \begin{minipage}[t]{\linewidth}\raggedright
Tái tạo container không dùng privileged mode.\strut
\end{minipage} & \begin{minipage}[t]{\linewidth}\raggedright
Bỏ privileged và thay bằng capability hoặc device cụ thể nếu bắt buộc.\strut
\end{minipage} \\
\hline
 & 5.6: Không mount thư mục hệ thống nhạy cảm của host & \begin{minipage}[t]{\linewidth}\raggedright
Loại bỏ bind mount nhạy cảm.\strut
\end{minipage} & \begin{minipage}[t]{\linewidth}\raggedright
Bỏ mount /, /etc, /proc, /sys, /dev; nếu bắt buộc thì read-only và có phê duyệt.\strut
\end{minipage} \\
\hline
 & 5.7: Không chạy sshd trong container & \begin{minipage}[t]{\linewidth}\raggedright
Loại bỏ SSH daemon khỏi image/container.\strut
\end{minipage} & \begin{minipage}[t]{\linewidth}\raggedright
Quản trị bằng docker exec, logs hoặc orchestration thay vì sshd trong container.\strut
\end{minipage} \\
\hline
 & 5.8: Không ánh xạ privileged port không cần thiết & \begin{minipage}[t]{\linewidth}\raggedright
Gỡ privileged port hoặc ghi nhận phê duyệt.\strut
\end{minipage} & \begin{minipage}[t]{\linewidth}\raggedright
Dùng port cao hơn hoặc reverse proxy; chỉ giữ port nhỏ hơn 1024 khi có nhu cầu rõ ràng.\strut
\end{minipage} \\
\hline
 & 5.9: Chỉ mở port cần thiết & \begin{minipage}[t]{\linewidth}\raggedright
Xóa các port publish không dùng.\strut
\end{minipage} & \begin{minipage}[t]{\linewidth}\raggedright
Cập nhật Docker CLI/Compose/Ansible để chỉ publish port nằm trong thiết kế.\strut
\end{minipage} \\
\hline
 & 5.10: Không chia sẻ network namespace của host & \begin{minipage}[t]{\linewidth}\raggedright
Recreate container trên network riêng.\strut
\end{minipage} & \begin{minipage}[t]{\linewidth}\raggedright
Bỏ network host; dùng bridge/overlay user-defined network.\strut
\end{minipage} \\
\hline
 & 5.11: Giới hạn memory cho container & \begin{minipage}[t]{\linewidth}\raggedright
Đặt memory limit theo workload.\strut
\end{minipage} & \begin{minipage}[t]{\linewidth}\raggedright
Dùng memory limit trong Docker CLI, playbook hoặc compose.\strut
\end{minipage} \\
\hline
 & 5.12: Đặt CPU priority phù hợp & \begin{minipage}[t]{\linewidth}\raggedright
Đặt CPU shares hoặc quota.\strut
\end{minipage} & \begin{minipage}[t]{\linewidth}\raggedright
Dùng cpu-shares hoặc CPU quota phù hợp, sau đó test hiệu năng.\strut
\end{minipage} \\
\hline
 & 5.13: Root filesystem ở chế độ read-only & \begin{minipage}[t]{\linewidth}\raggedright
Chạy container với rootfs read-only.\strut
\end{minipage} & \begin{minipage}[t]{\linewidth}\raggedright
Dùng read-only và khai báo tmpfs/volume riêng cho đường dẫn cần ghi.\strut
\end{minipage} \\
\hline
 & 5.14: Bind traffic vào interface cụ thể & \begin{minipage}[t]{\linewidth}\raggedright
Bind port vào IP được phê duyệt.\strut
\end{minipage} & \begin{minipage}[t]{\linewidth}\raggedright
Dùng 127.0.0.1:port:port hoặc IP NIC cụ thể thay vì 0.0.0.0.\strut
\end{minipage} \\
\hline
 & 5.15: Giới hạn restart policy & \begin{minipage}[t]{\linewidth}\raggedright
Đổi restart policy sang on-failure có giới hạn.\strut
\end{minipage} & \begin{minipage}[t]{\linewidth}\raggedright
Dùng restart on-failure:5 hoặc restart policy/retry tương ứng.\strut
\end{minipage} \\
\hline
 & 5.16: Không chia sẻ PID namespace của host & \begin{minipage}[t]{\linewidth}\raggedright
Recreate container không dùng host PID.\strut
\end{minipage} & \begin{minipage}[t]{\linewidth}\raggedright
Bỏ pid host khỏi cấu hình runtime.\strut
\end{minipage} \\
\hline
 & 5.17: Không chia sẻ IPC namespace của host & \begin{minipage}[t]{\linewidth}\raggedright
Recreate container không dùng host IPC.\strut
\end{minipage} & \begin{minipage}[t]{\linewidth}\raggedright
Bỏ ipc host; nếu cần chia sẻ IPC thì giới hạn giữa container được phê duyệt.\strut
\end{minipage} \\
\hline
 & 5.18: Không expose trực tiếp host device & \begin{minipage}[t]{\linewidth}\raggedright
Gỡ device mapping không cần thiết.\strut
\end{minipage} & \begin{minipage}[t]{\linewidth}\raggedright
Bỏ device mapping hoặc giới hạn đúng thiết bị/quyền truy cập theo nghiệp vụ.\strut
\end{minipage} \\
\hline
 & 5.19: Ghi đè ulimit khi cần & \begin{minipage}[t]{\linewidth}\raggedright
Khai báo ulimit rõ ràng.\strut
\end{minipage} & \begin{minipage}[t]{\linewidth}\raggedright
Dùng ulimit hoặc default-ulimits trong daemon.json đã được harden.\strut
\end{minipage} \\
\hline
 & 5.20: Không dùng shared mount propagation & \begin{minipage}[t]{\linewidth}\raggedright
Chuyển mount propagation sang private/slave.\strut
\end{minipage} & \begin{minipage}[t]{\linewidth}\raggedright
Cập nhật bind propagation và recreate container.\strut
\end{minipage} \\
\hline
 & 5.21: Không chia sẻ UTS namespace của host & \begin{minipage}[t]{\linewidth}\raggedright
Recreate container không dùng host UTS.\strut
\end{minipage} & \begin{minipage}[t]{\linewidth}\raggedright
Bỏ uts host khỏi cấu hình runtime.\strut
\end{minipage} \\
\hline
 & 5.22: Không tắt default seccomp profile & \begin{minipage}[t]{\linewidth}\raggedright
Bỏ cấu hình seccomp unconfined.\strut
\end{minipage} & \begin{minipage}[t]{\linewidth}\raggedright
Dùng default seccomp hoặc custom seccomp profile đã kiểm soát.\strut
\end{minipage} \\
\hline
 & 5.23: Không dùng docker exec với privileged option & \begin{minipage}[t]{\linewidth}\raggedright
Cấm quy trình vận hành dùng exec privileged.\strut
\end{minipage} & \begin{minipage}[t]{\linewidth}\raggedright
Cập nhật runbook, review quyền admin và bật logging/audit tập trung.\strut
\end{minipage} \\
\hline
 & 5.24: Không dùng docker exec với user root option & \begin{minipage}[t]{\linewidth}\raggedright
Vận hành bằng user ứng dụng.\strut
\end{minipage} & \begin{minipage}[t]{\linewidth}\raggedright
Dùng user ứng dụng hoặc sửa image để có USER non-root.\strut
\end{minipage} \\
\hline
 & 5.25: Xác nhận cgroup usage & \begin{minipage}[t]{\linewidth}\raggedright
Bỏ custom cgroup parent nếu không cần.\strut
\end{minipage} & \begin{minipage}[t]{\linewidth}\raggedright
Loại bỏ cgroup parent hoặc ghi tài liệu phê duyệt kèm resource control.\strut
\end{minipage} \\
\hline
 & 5.26: Ngăn container nhận thêm đặc quyền & \begin{minipage}[t]{\linewidth}\raggedright
Thêm no-new-privileges.\strut
\end{minipage} & \begin{minipage}[t]{\linewidth}\raggedright
Dùng security option no-new-privileges=true hoặc security opts trong Ansible.\strut
\end{minipage} \\
\hline
 & 5.27: Có runtime health check & \begin{minipage}[t]{\linewidth}\raggedright
Sửa healthcheck để phản ánh trạng thái thật.\strut
\end{minipage} & \begin{minipage}[t]{\linewidth}\raggedright
Thêm HEALTHCHECK hoặc healthcheck command nội bộ; nghiệm thu khi healthy.\strut
\end{minipage} \\
\hline
 & 5.28: Ghim phiên bản image, không dùng latest & \begin{minipage}[t]{\linewidth}\raggedright
Ghim image bằng version tag hoặc digest.\strut
\end{minipage} & \begin{minipage}[t]{\linewidth}\raggedright
Đổi nginx latest thành version tag hoặc sha256 digest đã xác minh.\strut
\end{minipage} \\
\hline
 & 5.29: Sử dụng PIDs cgroup limit & \begin{minipage}[t]{\linewidth}\raggedright
Đặt PID limit.\strut
\end{minipage} & \begin{minipage}[t]{\linewidth}\raggedright
Dùng pids limit trong Docker CLI, playbook hoặc compose.\strut
\end{minipage} \\
\hline
 & 5.30: Không dùng default bridge docker0 & \begin{minipage}[t]{\linewidth}\raggedright
Tạo và dùng network riêng.\strut
\end{minipage} & \begin{minipage}[t]{\linewidth}\raggedright
Chạy docker network create ten-network rồi recreate container trên network đó.\strut
\end{minipage} \\
\hline
 & 5.31: Không chia sẻ user namespace của host & \begin{minipage}[t]{\linewidth}\raggedright
Bỏ host user namespace.\strut
\end{minipage} & \begin{minipage}[t]{\linewidth}\raggedright
Bỏ userns host; cân nhắc userns-remap ở daemon nếu phù hợp.\strut
\end{minipage} \\
\hline
 & 5.32: Không mount Docker socket vào container & \begin{minipage}[t]{\linewidth}\raggedright
Gỡ Docker socket bind mount.\strut
\end{minipage} & \begin{minipage}[t]{\linewidth}\raggedright
Bỏ /var/run/docker.sock khỏi volumes; dùng API TLS hoặc công cụ không cần daemon socket.\strut
\end{minipage} \\
\hline
6 & Tránh Image sprawl & Xóa bỏ các images không sử dụng thủ công & Chạy lệnh docker image prune -a --force \\
\hline
& Tránh container sprawl & Xóa bỏ các containers đã dừng thủ công & Chạy lệnh docker container prune --force \\
\hline
\end{longtable}
}


Sau khi remediation thủ công hoàn tất, quản trị viên cần chạy lại playbook audit để xác nhận các cấu hình đã được khắc phục. Các thay đổi có ảnh hưởng đến Docker daemon, storage driver, network hoặc container đang chạy cần được thực hiện trong thời gian bảo trì để tránh làm gián đoạn dịch vụ.
